% ==============================================================================
% PRESENTASI SLIDE UAS ISR (BEAMER TEMPLATE)
% ==============================================================================

\documentclass[aspectratio=169]{beamer}

% Load Metadata (Relatif terhadap folder slides/)
% ==============================================================================
% METADATA MAKALAH UAS ISR
% ==============================================================================
% Silakan sesuaikan informasi di bawah ini dengan data kelompok Anda.

\newcommand{\JudulMakalah}{Perancangan dan Implementasi Sistem Robotika Berbasis Intelligent Systems}
\newcommand{\JudulSingkat}{Robotika Berbasis Intelligent Systems}
\newcommand{\MataKuliah}{Instrumentasi Sistem Robotika}
\newcommand{\TahunAkademik}{2026}
\newcommand{\DosenPengampu}{Dr.Eng. Dwi Joko Suroso, S.T., M.Eng.}
\newcommand{\ProgramStudi}{Departemen Teknik Nuklir dan Teknik Fisika}
\newcommand{\Fakultas}{Fakultas Teknik}
\newcommand{\Universitas}{Universitas Gadjah Mada}

% Anggota Kelompok
\newcommand{\AnggotaSatu}{Muhammad Khalis Farhan Azvi}
\newcommand{\NIMSatu}{22/493199/TK/54016}

\newcommand{\AnggotaDua}{Valen Fatli}
\newcommand{\NIMDua}{22/493755/TK/54109}

\newcommand{\AnggotaTiga}{Yasir Abdulaziz}
\newcommand{\NIMTiga}{22/497079/TK/54468}

\newcommand{\AnggotaEmpat}{Galih Salman Sadewo}
\newcommand{\NIMEmpat}{22/498121/TK/54625}

\newcommand{\AnggotaLima}{Faris Afif Ridyasmara}
\newcommand{\NIMLima}{22/502798/TK/54902}


% Tema Beamer (Clean and Professional)
\usetheme{Madrid}
\usecolortheme{default}

% Kustomisasi Warna Tema Madrid (Slate & Cyan)
\definecolor{primary}{RGB}{30, 41, 59}     % Slate 800 - Deep slate blue
\definecolor{secondary}{RGB}{14, 116, 144}   % Cyan 700 - Teal/Cyan
\definecolor{accent}{RGB}{249, 115, 22}      % Orange 500 - Orange accent

\setbeamercolor{palette primary}{bg=primary, fg=white}
\setbeamercolor{palette secondary}{bg=secondary, fg=white}
\setbeamercolor{palette tertiary}{bg=primary, fg=white}
\setbeamercolor{structure}{fg=secondary}
\setbeamercolor{titlelike}{parent=structure, fg=primary}
\setbeamercolor{block title}{bg=primary, fg=white}
\setbeamercolor{block body}{bg=black!5, fg=black}

% Pemaksaan Bahasa & Karakter
\usepackage[utf8]{inputenc}
\usepackage[indonesian]{babel}
\usepackage{graphicx}
\usepackage{booktabs}

% Lokasi Folder Gambar (Relatif terhadap folder slides/)
\graphicspath{{../figures/}}

% Detail Presentasi (Diambil otomatis dari config/metadata.tex)
\title[\JudulSingkat]{\texorpdfstring{\textbf{\JudulMakalah}}{\JudulMakalah}}
\subtitle{Tugas Akhir UAS \MataKuliah}
\author[Kelompok \AnggotaSatu{} dkk.]{%
  \texorpdfstring{%
    \textbf{Anggota Kelompok:}\\[0.2cm]
    \begin{tabular}{cc}
        \AnggotaSatu\ (\NIMSatu) & \AnggotaDua\ (\NIMDua) \\
        \AnggotaTiga\ (\NIMTiga) & \AnggotaEmpat\ (\NIMEmpat)
    \end{tabular}%
  }{Kelompok \AnggotaSatu, \AnggotaDua, \AnggotaTiga, \AnggotaEmpat}%
}
\institute[\Universitas]{%
  \texorpdfstring{%
    \Fakultas\\
    \Universitas%
  }{\Fakultas, \Universitas}%
}
\date{Semester Genap \TahunAkademik}

\begin{document}

% Slide 1: Judul & Anggota Kelompok (Sesuai Permintaan UAS)
\begin{frame}
    \titlepage
\end{frame}

% Slide 2: Pendahuluan (Latar Belakang & Masalah)
\begin{frame}{Pendahuluan}
    \begin{block}{Latar Belakang}
        Tuliskan ringkasan latar belakang proyek robotika kelompok Anda. Mengapa robot ini dirancang dan apa masalah nyata yang ingin diatasi.
    \end{block}
    \begin{block}{Rumusan Masalah}
        \begin{itemize}
            \item Bagaimana merancang arsitektur sistem robotika cerdas?
            \item Bagaimana performa sistem ketika diuji di lingkungan nyata?
        \end{itemize}
    \end{block}
\end{frame}

% Slide 3: Desain Perangkat Keras
\begin{frame}{Desain Perangkat Keras (Hardware)}
    \begin{columns}
        \begin{column}{0.5\textwidth}
            \textbf{Komponen Robot:}
            \begin{itemize}
                \item Mikrokontroler: Arduino Uno / ESP32
                \item Sensor Jarak: HC-SR04
                \item Driver Motor: L298N H-Bridge
                \item Aktuator: Motor DC Geared
            \end{itemize}
        \end{column}
        \begin{column}{0.5\textwidth}
            \centering
            % Gunakan gambar nyata robot kelompok Anda
            % \includegraphics[width=0.8\textwidth]{figures/foto_robot.jpg}
            \framebox[5cm]{\parbox[t][3cm]{4.5cm}{\vfill\centering [Foto Fisik Robot]\vfill}}
        \end{column}
    \end{columns}
\end{frame}

% Slide 4: Algoritma & Perangkat Lunak
\begin{frame}{Alur Logika Kendali Cerdas}
    \begin{itemize}
        \item Robot menggunakan algoritma \textit{Obstacle Avoidance} untuk mendeteksi rintangan di depan secara otonom.
        \item Pembacaan sensor jarak diproses setiap 100 milidetik.
        \item Keputusan gerakan (maju/belok kanan) diambil berdasarkan jarak ambang batas.
    \end{itemize}
    \begin{block}{Logika Navigasi}
        \texttt{IF Jarak < 20 cm THEN Belok\_Kanan() ELSE Maju\_Lurus()}
    \end{block}
\end{frame}

% Slide 5: Data Pengujian & Hasil
\begin{frame}{Data Pengujian Sensor Jarak}
    \centering
    \small
    \begin{tabular}{ccccc}
        \toprule
        \textbf{No.} & \textbf{Jarak Aktual (cm)} & \textbf{Jarak Sensor (cm)} & \textbf{Selisih / Error (cm)} & \textbf{Error (\%)} \\
        \midrule
        1  & 10.0 & 10.1 & 0.1 & 1.00 \\
        2  & 20.0 & 19.8 & 0.2 & 1.00 \\
        3  & 30.0 & 30.3 & 0.3 & 1.00 \\
        4  & 40.0 & 39.5 & 0.5 & 1.25 \\
        5  & 50.0 & 50.2 & 0.2 & 0.40 \\
        \midrule
        \multicolumn{4}{l}{\textbf{Rata-rata Error (\%)}} & \textbf{0.93\%} \\
        \bottomrule
    \end{tabular}
\end{frame}

% Slide 6: Kesimpulan & Saran
\begin{frame}{Kesimpulan \& Saran}
    \begin{block}{Kesimpulan}
        \begin{itemize}
            \item Integrasi sensor, aktuator, dan logika kendali berjalan sukses.
            \item Rata-rata error sensor ultrasonik sangat rendah yaitu $0.93\%$.
            \item Robot dapat menghindari rintangan dengan tingkat keberhasilan $90\%$.
        \end{itemize}
    \end{block}
    
    \begin{block}{Saran Pengembangan}
        Menambahkan sensor inframerah tambahan untuk navigasi sisi kiri-kanan serta menerapkan algoritma kendali PID/Fuzzy.
    \end{block}
\end{frame}

% Slide 7: Penutup
\begin{frame}
    \centering
    \Huge \textbf{Terima Kasih}\\[0.5cm]
    \large Ada Pertanyaan?
\end{frame}

\end{document}
